\begin{tabular}{lll}
\toprule
OEIS & entry name & closed form in our $N$ \\
\midrule
\rt{A023434} & \parbox[t]{0.44\textwidth}{\raggedright Dying rabbits: a(n) = a(n-1) + a(n-2) - a(n-4).} & \parbox[t]{0.25\textwidth}{\raggedright a(N) = a(N-1) + a(N-2) - a(N-4)} \\
\rt{A000931} & \parbox[t]{0.44\textwidth}{\raggedright Padovan sequence (or Padovan numbers): a(n) = a(n-2) + a(n-3) with a(0) = 1, a(1) = a(2) = 0.} & \parbox[t]{0.25\textwidth}{\raggedright a(N) = a(N-2) + a(N-3)} \\
\rt{A000930} & \parbox[t]{0.44\textwidth}{\raggedright Narayana's cows sequence: a(0) = a(1) = a(2) = 1; thereafter a(n) = a(n-1) + a(n-3).} & \parbox[t]{0.25\textwidth}{\raggedright a(N) = a(N-1) + a(N-3)} \\
\rt{A179070} & \parbox[t]{0.44\textwidth}{\raggedright a(1)=a(2)=a(3)=1, a(4)=3; thereafter a(n) = a(n-1) + a(n-3).} & \parbox[t]{0.25\textwidth}{\raggedright a(N) = a(N-1) + a(N-3)} \\
\rt{A173862} & \parbox[t]{0.44\textwidth}{\raggedright a(n) = A158772(n-1)/21.} & \parbox[t]{0.25\textwidth}{\raggedright 2\textasciicircum{}floor((N+1)/3)} \\
\rt{A000045} & \parbox[t]{0.44\textwidth}{\raggedright Fibonacci numbers: F(n) = F(n-1) + F(n-2) with F(0) = 0 and F(1) = 1.} & \parbox[t]{0.25\textwidth}{\raggedright F(N+2)} \\
\rt{A000045} & \parbox[t]{0.44\textwidth}{\raggedright Fibonacci numbers: F(n) = F(n-1) + F(n-2) with F(0) = 0 and F(1) = 1.} & \parbox[t]{0.25\textwidth}{\raggedright F(N)} \\
\rt{A001609} & \parbox[t]{0.44\textwidth}{\raggedright a(1) = a(2) = 1, a(3) = 4; thereafter a(n) = a(n-1) + a(n-3).} & \parbox[t]{0.25\textwidth}{\raggedright a(N) = a(N-1) + a(N-3)} \\
\rt{A169985} & \parbox[t]{0.44\textwidth}{\raggedright Round phi\textasciicircum{}n to the nearest integer.} & \parbox[t]{0.25\textwidth}{\raggedright round(phi\textasciicircum{}N)} \\
\rt{A005251} & \parbox[t]{0.44\textwidth}{\raggedright a(0) = 0, a(1) = a(2) = a(3) = 1; thereafter, a(n) = a(n-1) + a(n-2) + a(n-4).} & \parbox[t]{0.25\textwidth}{\raggedright a(N) = a(N-1) + a(N-2) + a(N-4)} \\
\rt{A005251} & \parbox[t]{0.44\textwidth}{\raggedright a(0) = 0, a(1) = a(2) = a(3) = 1; thereafter, a(n) = a(n-1) + a(n-2) + a(n-4).} & \parbox[t]{0.25\textwidth}{\raggedright a(N) = a(N-1) + a(N-2) + a(N-4)} \\
\rt{A000045} & \parbox[t]{0.44\textwidth}{\raggedright Fibonacci numbers: F(n) = F(n-1) + F(n-2) with F(0) = 0 and F(1) = 1.} & \parbox[t]{0.25\textwidth}{\raggedright F(N+2)} \\
\rt{A008619} & \parbox[t]{0.44\textwidth}{\raggedright Positive integers repeated.} & \parbox[t]{0.25\textwidth}{\raggedright floor((N+1)/2) + 1} \\
\rt{A004525} & \parbox[t]{0.44\textwidth}{\raggedright One even followed by three odd.} & \parbox[t]{0.25\textwidth}{\raggedright (N + 1 - eps)/2 + 1, eps = ceil(N/2) mod 2} \\
\rt{A000045} & \parbox[t]{0.44\textwidth}{\raggedright Fibonacci numbers: F(n) = F(n-1) + F(n-2) with F(0) = 0 and F(1) = 1.} & \parbox[t]{0.25\textwidth}{\raggedright F(N)} \\
\rt{A000045} & \parbox[t]{0.44\textwidth}{\raggedright Fibonacci numbers: F(n) = F(n-1) + F(n-2) with F(0) = 0 and F(1) = 1.} & \parbox[t]{0.25\textwidth}{\raggedright F(N+2)} \\
\rt{A178715} & \parbox[t]{0.44\textwidth}{\raggedright a(n) = solution to the "Select All, Copy, Paste" problem: Given the ability to type a single letter, or to type individual "Select All", "Copy" or "Paste" command keystrokes, what is the maximal number of letters of text that can be obtained with n keystrokes?} & \parbox[t]{0.25\textwidth}{\raggedright a(N) = 4 a(N-5) for N >= 16} \\
\rt{A214282} & \parbox[t]{0.44\textwidth}{\raggedright Largest Euler characteristic of a downset on an n-dimensional cube.} & \parbox[t]{0.25\textwidth}{\raggedright ---} \\
\rt{A001405} & \parbox[t]{0.44\textwidth}{\raggedright a(n) = binomial(n, floor(n/2)).} & \parbox[t]{0.25\textwidth}{\raggedright binomial(N+1, floor((N+1)/2))} \\
\rt{A006498} & \parbox[t]{0.44\textwidth}{\raggedright a(n) = a(n-1) + a(n-3) + a(n-4), a(0) = a(1) = a(2) = 1, a(3) = 2.} & \parbox[t]{0.25\textwidth}{\raggedright F(floor(N/2)+2) * F(ceil(N/2)+2)} \\
\rt{A195971} & \parbox[t]{0.44\textwidth}{\raggedright Number of n X 1 0..4 arrays with each element x equal to the number its horizontal and vertical neighbors equal to 2,0,1,3,4 for x=0,1,2,3,4.} & \parbox[t]{0.25\textwidth}{\raggedright F(m+1)\textasciicircum{}2 (N=2m even), F(m) F(m+3) (N=2m+1 odd)} \\
\rt{A008619} & \parbox[t]{0.44\textwidth}{\raggedright Positive integers repeated.} & \parbox[t]{0.25\textwidth}{\raggedright floor(N/2) + 1} \\
\rt{A000034} & \parbox[t]{0.44\textwidth}{\raggedright Period 2: repeat [1, 2]; a(n) = 1 + (n mod 2).} & \parbox[t]{0.25\textwidth}{\raggedright 1 + (N mod 2)} \\
\rt{A007877} & \parbox[t]{0.44\textwidth}{\raggedright Period 4 zigzag sequence: repeat [0,1,2,1].} & \parbox[t]{0.25\textwidth}{\raggedright period 4: 1, 0, 1, 2 for N = 1, 2, 3, 0 (mod 4)} \\
\bottomrule
\end{tabular}
